\documentclass{ip-quiz}

\quizlevel{n2}
\quiznum{2}
\quiztitle{Strings y condicionales}

\begin{document}

\makeheader
\maketitle

\begin{sectionbox}{Enunciado}
Escriba una función que reciba una lista de números y retorne algún elemento que no esté repetido en la lista. Si no hay ningún elemento que cumpla esta condición, retorne \texttt{-1}.
\end{sectionbox}

\begin{sectionbox}{Solución}
\begin{minted}[escapeinside=||]{python}
def elemento_no_repetido(numeros: list) -> int:
    conteos = {}
    |\tms{forn}|for numero in numeros:|\tme{forn}|
        conteos[numero] = conteos.get(numero, 0) + 1

    i = 0
    resultado = -1
    while i < len(numeros) and |\tms{cent}|resultado == -1:|\tme{cent}|
        if conteos[numeros[i]] == 1:
            resultado = numeros[i]
        i += 1
    return resultado
\end{minted}

\begin{tikzpicture}[remember picture, overlay]

  \node[box, fit=(forns)(forne)] (cforn) {};
  \node[note, anchor=north west, text width=5cm]
    (noteForn) at ([xshift=-12cm, yshift=0.8cm] current page.east |- forns)
    {El \texttt{for} sobre un diccionario recorre solo las llaves.};
  \draw[arr, dotted] (cforn.east) to[bend right=5] (noteForn.west);

  \node[box, fit=(cents)(cente)] (ccent) {};
  \node[note, anchor=north west, text width=5cm]
    (noteCent) at ([xshift=-9cm, yshift=0.8cm] current page.east |- cents)
    {Este es un centinela sin necesidad de ser booleano: su valor inicial \texttt{-1} indica que aún no se encontró resultado.};
  \draw[arr, dotted] (ccent.east) to[bend right=5] (noteCent.west);

\end{tikzpicture}
\end{sectionbox}

Notas:
\begin{itemize}
    \item Hay varias otras soluciones a este ejercicio. Algunas utilizan métodos como \texttt{count} o \texttt{replace} en lugar de hacer dos ciclos consecutivos. Sin embargo, llamar \texttt{count} o \texttt{replace} dentro de un ciclo obliga a recorrer toda la cadena en cada iteración del ciclo, lo cual hace la solución más lenta. Son correctas, pero menos eficientes.
    \item El segundo ciclo debe ser un recorrido parcial, puesto que se está buscando un elemento con una característica específica. Utilizar \texttt{for} o hacerlo como un recorrido total de cualquier otra manera es incorrecto.
\end{itemize}





\end{document}
